\documentclass[english, parskip=full]{scrartcl}
\usepackage[utf8]{inputenc}  % Kodierung der Datei
\usepackage[english]{babel}  % Sprache auf Deutsch 
\usepackage{color}
\usepackage{graphicx, gensymb}
\usepackage{amsfonts, cancel, amssymb}
\usepackage{amsmath, amsthm, array, esint}
\usepackage{booktabs, multirow}  % schönere Tabellen
\usepackage{caption}  % Anpassung der Captions
\usepackage{scrlayer-scrpage}    % Manipulation des Seitenstils
\pagestyle{scrheadings}  % Stil für die Seite setzen
\automark{section}  % Abschnittsnamen als Seitenbeschriftung 
\setheadsepline{.5pt}    % Kopzeile durch Linie abtrennen
\usepackage[hidelinks]{hyperref}  % Links und weitere PDF-Features
\usepackage{typearea, tikz, pgffor, verbatim}
\usetikzlibrary{calc, quotes}
\usepackage[cache=false, outputdir=build]{minted}
\usepackage{placeins} % provides \FloatBarrier

\definecolor{bg}{rgb}{0.9,0.9,0.9}
\newcommand\numberthis{\addtocounter{equation}{1}\tag{\theequation}}
\newcommand{\bra}{}
\newcommand{\kom}{}
\newcommand{\brak}{}
\newcommand{\braket}{}
\newcommand{\intx}{}
\newcommand{\intr}{}
\newcommand{\kla}{}
\newcommand{\klam}{}
\newcommand{\klammer}{}
\newcommand{\oper}{}
\newcommand{\tecxt}{}
\newcommand{\Bra}{}
\newcommand{\Ket}{}
\renewcommand{\i}{\text{i}}
\newcommand{\e}{\text{e}}
\newcommand*{\Fourier}{\textsc{Fourier}\ }
\newcommand*{\F}{\mathcal{F}}
\newcommand*{\sinc}{\text{sinc\,}}
\newcommand{\conv}{\mathop{\scalebox{3.5}{\raisebox{-0.38ex}{\makebox[0.5\width][c]{$\ast$}}}}\limits}

\newcommand*{\titel}{03 Properties of Wannier functions}
\newcommand*{\autor}{Joachim Schwardt}

\hypersetup{pdfauthor={\autor}, pdftitle={\titel}}

%\titlehead{titlehead}
\subject{Topological Insulators and Superconductors}
\title{\titel}
\author{\autor}
\date{\today}

\begin{document}

%\begin{titlepage}
\maketitle  % Titel setzen
%\tableofcontents  % Inhaltsverzeichnis setzen
%\end{titlepage}

\section*{Problem 3.1: Orthonormality of Wannier functions}
We are given the relation
\begin{align}
\brak\langle \psi_\mathbf{k} | \chi_{\mathbf{k}'} \rangle &= \frac{V_\tecxt\text{cell}}{(2\pi)^3} \brak\langle u_\mathbf{k} | v_{\mathbf{k}'} \rangle \delta\kla\left( \mathbf{k} - \mathbf{k}' \right).
\end{align}
The Wannier functions are defined as
\begin{align}
\Ket | \omega_{n \mathbf{R}} \rangle &:= \frac{V_\tecxt\text{cell}}{(2\pi)^3} \intx\int_{\tecxt\text{BZ}}^{ } \e^{-\i\mathbf{k}\cdot\mathbf{R}} | \psi_{n\mathbf{k}} \rangle\, \mathrm{d}^3\mathbf{k}.
\end{align}
Note that $\psi_{n\mathbf{k}}(\mathbf{R}) = \e^{\i\mathbf{k}\cdot \mathbf{R}} u_{n\mathbf{k}} (\mathbf{R})$ and $\chi_{n\mathbf{k}}(\mathbf{R}) = \e^{\i\mathbf{k}\cdot \mathbf{R}} v_{n\mathbf{k}} (\mathbf{R})$.
Thus we find
\begin{align}
\brak\langle \omega_{n\mathbf{R}} | \omega_{n'\mathbf{R}'} \rangle &= \frac{V_\tecxt\text{cell}^2}{(2\pi)^6} \intx\int_{\tecxt\text{BZ}^2}^{ } \e^{\i\mathbf{k}\cdot\mathbf{R}} \e^{-\i\mathbf{k}'\cdot\mathbf{R}'} \brak\langle \psi_{n\mathbf{k}}  | \psi_{n'\mathbf{k}'} \rangle  \, \mathrm{d}^3\mathbf{k}\, \tecxt\text{d}^3 \mathbf{k}' = \\
%&= \delta_{nn'}  \frac{V_\tecxt\text{cell}^2}{(2\pi)^6} \intx\int_{\tecxt\text{BZ}^2}^{ } \e^{\i\mathbf{k}\cdot\mathbf{R}} \e^{-\i\mathbf{k}'\cdot\mathbf{R}'} \intr\int_{V_\tecxt\text{cell}} \brak\langle \psi_{\mathbf{k}}  |   \mathbf{r} \rangle\langle \mathbf{r}   | \psi_{\mathbf{k}'} \rangle \, \mathrm{d}^{3}\mathbf{r} \, \mathrm{d}^3\mathbf{k}\, \tecxt\text{d}^3 \mathbf{k}' = \\
%&= \delta_{nn'}  \frac{V_\tecxt\text{cell}^2}{(2\pi)^6} \intr\int_{V_\tecxt\text{cell}} \intx\int_{\tecxt\text{BZ}^2}^{ } \e^{ \i\mathbf{k}\cdot\mathbf{R}} \e^{-\i\mathbf{k}'\cdot\mathbf{R}'}  u_\mathbf{k} (\mathbf{r}) u_{\mathbf{k}'}^\ast (\mathbf{r}) \e^{\i (\mathbf{k}' - \mathbf{k})\cdot\mathbf{r}}  \, \mathrm{d}^3\mathbf{k}\, \tecxt\text{d}^3 \mathbf{k}'\, \mathrm{d}^{3}\mathbf{r} = \\
%&= \delta_{nn'} \intx\int_{V_\tecxt\text{cell}}^{ } \delta (\mathbf{r} - \mathbf{R}) \delta (\mathbf{r} - \mathbf{R}') u_\mathbf{k} (\mathbf{r}) u_{\mathbf{k}'}^\ast (\mathbf{r}) \, \mathrm{d}^3\mathbf{r} = \\
%&= \frac{V_\tecxt\text{cell}^2}{(2\pi)^6} \intx\int_{\tecxt\text{BZ}^2}^{ } \e^{\i\mathbf{k}\cdot\mathbf{R}} \e^{-\i\mathbf{k}'\cdot\mathbf{R}'} \brak\langle \psi_{n\mathbf{k}}  | \psi_{n'\mathbf{k}'} \rangle  \, \mathrm{d}^3\mathbf{k}\, \tecxt\text{d}^3 \mathbf{k}' = \\
&=  \delta_{nn'} \frac{V_\tecxt\text{cell}}{(2\pi)^3} \intx\int_{\tecxt\text{BZ}^2}^{ } \e^{\i\mathbf{k}\cdot\mathbf{R}} \e^{-\i\mathbf{k}'\cdot\mathbf{R}'} \delta \kla\left( \mathbf{k} - \mathbf{k}' \right)  \, \mathrm{d}^3\mathbf{k}\, \tecxt\text{d}^3 \mathbf{k}' = \\
&=  \delta_{nn'} \frac{V_\tecxt\text{cell}}{(2\pi)^3} \intx\int_{\tecxt\text{BZ}}^{ } \e^{\i\mathbf{k} \cdot \kla\left( \mathbf{R} - \mathbf{R}' \right)} \, \mathrm{d}^3\mathbf{k} = \\
&= \delta_{nn'}\delta \kla\left( \mathbf{R} - \mathbf{R}' \right).
\end{align}
Here we only made use of the identity
\begin{align}
\brak\langle \psi_{n\mathbf{k}}  | \psi_{n'\mathbf{k}'} \rangle &= \delta_{nn'} \intx\int_{V_\tecxt\text{cell}}^{ } \e^{-\i\kla\left( \mathbf{k} - \mathbf{k}' \right) \cdot\mathbf{r}}  u_\mathbf{k} (\mathbf{r}) u_{\mathbf{k}'}^\ast (\mathbf{r}) \, \mathrm{d}^3\mathbf{r} = \delta_{nn'} \frac{(2\pi)^3}{V_\tecxt\text{cell}} \delta\kla\left( \mathbf{k} - \mathbf{k}' \right).
\end{align}

\newpage
\section*{Problem 3.2: }
Consider the quantity
\begin{align}
\Omega &:= \sum_n^\tecxt\text{occ} \brak\langle \partial_\lambda u_{n\mathbf{k}} | \partial_\mathbf{k} u_{n\mathbf{k}} \rangle .
\end{align}
Note that $\brak\langle u_{n\mathbf{k}} | u_{n\mathbf{k}} \rangle = 1$ and therefore
\begin{align}
\partial \brak\langle u_{n\mathbf{k}} | u_{n\mathbf{k}} \rangle &= \brak\langle u_{n\mathbf{k}} | \partial u_{n\mathbf{k}} \rangle + \brak\langle \partial u_{n\mathbf{k}} | u_{n\mathbf{k}} \rangle = 0
\end{align}
for any differential $\partial$.
Furthermore, we find
\begin{align}
\partial_\lambda \brak\langle u_{n\mathbf{k}} | \partial_\mathbf{k} u_{n\mathbf{k}} \rangle &= \brak\langle \partial_\lambda u_{n\mathbf{k}} | \partial_\mathbf{k} u_{n\mathbf{k}} \rangle + \brak\langle  u_{n\mathbf{k}} | \partial_\lambda \partial_\mathbf{k} u_{n\mathbf{k}} \rangle ,\\
\partial_\lambda \brak\langle \partial_\mathbf{k} u_{n\mathbf{k}} |  u_{n\mathbf{k}} \rangle &= \brak\langle \partial_\lambda \partial_\mathbf{k} u_{n\mathbf{k}} |  u_{n\mathbf{k}} \rangle + \brak\langle \partial_\mathbf{k} u_{n\mathbf{k}} | \partial_\lambda  u_{n\mathbf{k}} \rangle.
\end{align}
The sum of these expression is zero according to the above identity, which leaves us with the sum of the real parts vanishing, i.e.
\begin{align}
0 &= \tecxt\text{Re\,} \brak\langle \partial_\lambda \partial_\mathbf{k} u_{n\mathbf{k}} |  u_{n\mathbf{k}} \rangle + \tecxt\text{Re\,} \brak\langle \partial_\mathbf{k} u_{n\mathbf{k}} | \partial_\lambda  u_{n\mathbf{k}} \rangle.
\end{align}
We can make use of this by considering the real part of $\Omega$, since 
\begin{align}
\tecxt\text{Re\,} \Omega &= -\sum_{n}^{\tecxt\text{occ}} \tecxt\text{Re\,} \brak\langle u_{n\mathbf{k}} |  \partial_\lambda \partial_\mathbf{k} u_{n\mathbf{k}} \rangle
\end{align}
This is not really going anywhere. 
I do not have an idea to make any real progress here unfortunately.
All I know is that this expression showed up in the final equation of the lecture, but there we asked for the imaginary part of the solution. 
If this expression were fully real, the lecture result would thus be zero. 
Therefore something is wrong ...

\newpage
\section*{Problem 3.3: Integral transformation for the Berry curvature}
Let $\mathbf{k} = (k_x,k_y)^\top$, $\Omega_{\mu\nu} = -2\tecxt\text{Im\,} \brak\langle \partial_\mu u_{n\mathbf{k}} | \partial_\nu u_{n\mathbf{k}} \rangle = \epsilon_{\mu\nu}\Omega$ and $\kappa = (\kappa_1, \kappa_2)^\top$, $\bar{\Omega}_{jm}= -2\tecxt\text{Im\,} \brak\langle \partial_j u_{n\kappa} | \partial_m u_{n\kappa} \rangle = \epsilon_{jm} \bar{\Omega}$.
Assume that the two frames are related by $k_\mu = B_{\mu j}\kappa_j$.
Then the multi-dimensional substitution rule for integrals gives
\begin{align}
\tecxt\text{d}^2\mathbf{k} &= |J| \tecxt\text{d}^2\vec{\kappa},
\end{align}
where the Jacobian is
\begin{align}
J &= \det J_{\mu\nu} = \det \kla\left( \frac{\partial k_\mu }{\partial \kappa_j} \right) = \det \kla\left( B_{\mu j} \right) = \det B.
\end{align}
Note that in more explicit notation, we have
\begin{align}
\epsilon_{jm} B_{\mu j}B_{\nu m} &= B_{\mu 1}B_{\nu 2} - B_{\mu 2}B_{\nu 1}.
\end{align}
For $\mu=\nu$, this obviously gives
\begin{align}
\epsilon_{jm} B_{\mu j}B_{\nu m} &= B_{\mu 1}B_{\mu 2} - B_{\mu 2}B_{\mu 1} = 0.
\end{align}
For $\mu = 1$ and $\nu = 2$ we instead have
\begin{align}
\epsilon_{jm} B_{\mu j}B_{\nu m} &=  B_{1 1}B_{2 2} - B_{1 2}B_{2 1} = \det B.
\end{align}
For $\mu = 2$ and $\nu = 1$ we find 
\begin{align}
\epsilon_{jm} B_{\mu j}B_{\nu m} &=  B_{2 1}B_{1 2} - B_{2 2}B_{1 1} = -\det B,
\end{align}
which allows us to summarize the results as
\begin{align}
\epsilon_{jm} B_{\mu j}B_{\nu m} &=  \epsilon_{\mu\nu}\det B.
\end{align}
Then using $\partial_j = B_{\mu j}\partial_\mu$ we find
\begin{align}
\bar{\Omega}_{jm} &= -2\tecxt\text{Im\,}\brak\langle \partial_j u_{n\kappa} | \partial_m u_{n\kappa} \rangle = -2\tecxt\text{Im\,}\brak\langle B_{\mu j}\partial_\mu u_{n\mathbf{k}} | B_{\nu m}\partial_\nu u_{n\mathbf{k}} \rangle = \\
%&= -2\tecxt\text{Im\,} B_{\mu j}B_{\nu m} \brak\langle \partial_\mu u_{n\mathbf{k}} | \partial_\nu u_{n\mathbf{k}} \rangle  = \\
&= B_{\mu j}B_{\nu m} \Omega_{\mu\nu} = B_{\mu j}B_{\nu m} \epsilon_{\mu\nu} \Omega = \kla\left( B_{1j}B_{2m} - B_{2j}B_{1m} \right) \Omega = \\
&=  \det (B) \epsilon_{jm} \Omega.
\end{align}
Now we can show that the following integral is invariant under the described transformation, since
\begin{align}
\intx\int_{\tecxt\text{BZ}}^{ } \Omega \, \mathrm{d}^2\mathbf{k} &= \intx\int_{\tecxt\text{BZ}}^{ } \Omega_{12} \, \det
 (B) \mathrm{d}^2\kappa = \intx\int_{\tecxt\text{BZ}}^{ } \bar{\Omega}_{12} \, \mathrm{d}^2\kappa = \intx\int_{\tecxt\text{BZ}}^{ } \bar{\Omega} \, \mathrm{d}^2\kappa.
\end{align}


%\begin{thebibliography}{9}
%\bibitem{example} description, \url{https://www.exmapleurl}, day.\,month~2021.
%\end{thebibliography}

\end{document}